\uuid{vfCR}
\exo7id{1822}
\titre{exo7 1822}
\auteur{drutu}
\organisation{exo7}
\datecreate{2003-10-01}
\isIndication{false}
\isCorrection{false}
\chapitre{Fonction de plusieurs variables}
\sousChapitre{Différentiabilité}
\cpge{mpsi-mp2i.fonctions-deux-variables.b, pcsi.fonctions-deux-variables.b, ptsi.fonctions-deux-variables.b, bcpst1.fonctions-deux-variables.c, mp-mpi.calcul-differentiel-optimisation.a, mp-mpi.calcul-differentiel-optimisation.b, pc.fonctions-plusieurs-variables.c, psi.fonctions-plusieurs-variables.c}
\module{Analyse}
\niveau{L2}
\difficulte{}

\contenu{
\question{
Calculer la d\'eriv\'ee de la fonction $F(x,y,z)=\ln \left( 
e^x+e^y+e^z\right)$ \`a l'origine dans une direction formant avec les axes de 
coordonn\'ees $x,y,z$ les angles $\alpha ,\beta , \gamma$.
}
}
